\begin{table*}[!t]
\centering
\caption{Baseline provenance and local adaptation in the common evaluation pipeline.}
\label{tab:baselines}
\scriptsize
\renewcommand{\arraystretch}{1.08}
\begin{tabularx}{\textwidth}{@{}p{0.10\textwidth}p{0.18\textwidth}p{0.21\textwidth}X p{0.12\textwidth}@{}}
\toprule
Method & Source family & Original objective & Local adaptation and tuning protocol & Fidelity \\
\midrule
DDQN & DQN/DDQN/dueling value learning \cite{ref35,ref36,ref41} & Discrete action-value control with target-network stabilization. & Centralized binary edge/cloud policy trained on the same 30 features, purged split, validation checkpointing, and test-only metrics. & Faithful DDQN-style baseline. \\
MTOSA & Classical MEC offloading and scheduling family \cite{ref17,ref24,ref28} & Task placement under latency, computation, and communication constraints. & Implemented as a deterministic scheduling comparator using the same latency, rejection, and scenario scoring functions. & Style-inspired. \\
FL-DDPG & DDPG actor-critic plus federated MEC learning family \cite{ref02,ref04,ref05,ref38} & Continuous-control actor-critic learning under distributed data. & Adapted as a federated actor-critic comparator inside the same train/validation/test protocol and binary destination evaluation. & Style-inspired comparator. \\
GTPSO & Population/metaheuristic offloading family \cite{ref28} & Search over destination choices using swarm-like candidate updates. & Implemented as a PSO-inspired policy comparator scored with the shared held-out evaluation functions. & Style-inspired. \\
PTS-RA & Priority and resource-aware scheduling family \cite{ref18,ref19} & Allocate computation and communication resources under QoS and energy constraints. & Implemented as a priority/resource heuristic over the same task, edge, channel, and cloud state features. & Style-inspired. \\
JTOS & Joint task-offloading and scheduling family \cite{ref03,ref07,ref12} & Joint destination and scheduling decisions for MEC task streams. & Adapted to the binary destination setting while preserving the common split, feature, and metric definitions. & Style-inspired. \\
\fedddqn{} & Proposed federated value-learning design \cite{ref30,ref31,ref32,ref36,ref41} & Zone-aware binary offloading with QoS-aligned reward and validation-selected checkpointing. & Main trainable contribution; validation selects the checkpoint, and the final test interval is used only once for reporting. & Proposed method. \\
Oracle & Generated latency-reference decision rule & Non-trainable rule that selects the lower finite generated latency among edge and cloud alternatives. & Used only to estimate remaining latency headroom; not available to deployable policies and not a universal QoS ceiling. & Latency reference only. \\
\bottomrule
\end{tabularx}
\end{table*}
